\documentclass{report}
\usepackage[utf8]{inputenc}
\usepackage[spanish]{babel}
\usepackage{amsmath}
\usepackage{amssymb}
\usepackage{amsthm}

\theoremstyle{definition}
\newtheorem{definicion}{Definición}[chapter]
\theoremstyle{plain}
\newtheorem{teorema}{Teorema}[chapter]
\newtheorem{lema}{Lema}[chapter]
\theoremstyle{remark}
\newtheorem{ejemplo}{Ejemplo}[chapter]

\begin{document}

\setcounter{chapter}{1}
\chapter{Gramáticas de concatenación de rango y el SAT como lenguaje formal}
\label{cap:rcg-sat}

En este capítulo se presentan las gramáticas de concatenación de rango,
formalismo de base sobre el que se desarrolla el resto del trabajo. Se
analizan sus principales propiedades y limitaciones, y se revisan los
antecedentes directos que abordan el problema de la satisfacibilidad
booleana mediante formalismos de la teoría de lenguajes formales, con
énfasis en los trabajos previos realizados en la Facultad de Matemática
y Computación de la Universidad de La Habana. Los conceptos básicos
utilizados en este capítulo se presentaron en el capítulo anterior.

El capítulo se organiza en tres secciones. La sección~\ref{sec:rcg}
presenta las gramáticas de concatenación de rango, su presentación
intuitiva, sus definiciones formales y un ejemplo completo de
reconocimiento. La sección~\ref{sec:propiedades-rcg} expone las
propiedades de clausura del formalismo y su caracterización en
términos de complejidad. La sección~\ref{sec:antecedentes} revisa los
antecedentes directos del trabajo: los estudios previos en la Facultad
que abordan el SAT mediante formalismos de la teoría de lenguajes.

\section{Gramáticas de concatenación de rango}
\label{sec:rcg}

Las gramáticas de concatenación de rango (RCG)~\cite{boullier1998proposal}
son un formalismo propuesto en 1998 por Pierre Boullier, un investigador
en el campo de la lingüística computacional. El objetivo inicial del
formalismo era proporcionar un modelo más expresivo que las gramáticas
libres del contexto para describir ciertos fenómenos del lenguaje
natural, como los números chinos y el orden variable de algunas palabras
en alemán~\cite{boullier1999chinese}.

Antes de entrar a la presentación del formalismo, es necesario precisar
una convención terminológica que rige a lo largo de todo el documento.
El formalismo que se presenta en este capítulo corresponde, en la
literatura, a las \textbf{gramáticas de concatenación de rango
positivas} (PRCG), una de las dos variantes del formalismo general de
Boullier. La otra variante, las gramáticas de concatenación de rango con
negación, admite cláusulas con predicados negados y no se utiliza en
este trabajo. La restricción a la variante positiva es importante
porque, como se establece más adelante, la clase de lenguajes
reconocidos por una PRCG coincide exactamente con la clase $PTIME$, y
los resultados de los capítulos posteriores se apoyan en dicha
caracterización. Para abreviar, a lo largo del documento se adopta el
nombre \textbf{gramática de concatenación de rango} (RCG) para
referirse, salvo mención explícita, a la variante positiva.

La sección se organiza en tres partes. Primero se introducen de manera
intuitiva los elementos de las gramáticas de concatenación de rango.
Después se presentan las definiciones formales. Por último se presenta
un ejemplo completo de reconocimiento por una RCG.

\subsection{Presentación intuitiva}

En esta subsección se presentan las nociones de predicado, argumento,
sustitución de rango y derivación, que sirven de base para las
definiciones formales de la siguiente subsección.

A los no terminales de una RCG se les llama \textbf{predicados}. Cada
predicado tiene un conjunto de argumentos. A la cantidad de argumentos
de un predicado se le denomina \textbf{aridad}. Por ejemplo, $A(X, Y)$
denota el predicado $A$ con argumentos $X$ e $Y$; en este caso, la
aridad de $A$ es $2$.

Cada argumento de un predicado puede estar formado por variables y
terminales. Por convenio, las variables se denotan por letras mayúsculas
del final del alfabeto, y los terminales con letras minúsculas. Con este
convenio, el predicado $B(aX, XY, abZ)$ tiene aridad $3$. El primer
argumento está formado por el terminal $a$ y la variable $X$. El
segundo argumento está formado por la concatenación de las variables
$X$ e $Y$. El tercer argumento está formado por la concatenación de los
terminales $a$ y $b$ y la variable $Z$.

Cada predicado reconoce un vector de cadenas cuya dimensión es la
aridad del predicado. Cada cadena del vector se asocia a un argumento
del predicado.

Por ejemplo, si al predicado $A(X, Y)$ se le asocia el vector $[abc,
d]$, al primer argumento de $A$ se le asigna la cadena $abc$ y al
segundo la cadena $d$. Esto significa que $X = abc$ y que $Y = d$.

De manera similar, si al predicado $B(aX, XY, abZ)$ se le asigna el
vector $[aa, de, abcc]$, al primer argumento de $B$ se le asigna la
cadena $aa$, al segundo la cadena $de$ y al tercero la cadena $abcc$.
Es decir, $aX = aa$, $XY = de$ y $abZ = abcc$. Como $X$ es una variable
y se tiene que $aX = aa$, el único valor que puede tomar $X$ es $X = a$.
En el tercer argumento ocurre algo similar: como $abZ = abcc$, para que
se cumpla la igualdad el único valor posible para $Z$ es $Z = cc$. En el
segundo argumento, las variables $X$ e $Y$ deben tomar valores tales que
su concatenación sea $de$; esto se puede hacer de tres formas:
$X = d$, $Y = e$; $X = de$, $Y = \varepsilon$; y $X = \varepsilon$,
$Y = de$.

Si al predicado $B(aX, XY, abZ)$ se le asigna el vector
$[bb, de, abcc]$, entonces al tenerse $aX = bb$, $X$ no puede tomar
ningún valor: el primer carácter $b$ de la cadena no coincide con el
terminal $a$ que encabeza el argumento $aX$.

A las producciones de esta gramática se les denomina \textbf{cláusulas}.
La parte izquierda de la producción siempre está formada por un único
predicado. La parte derecha puede contener cualquier cantidad de
predicados con sus respectivos argumentos, o la cadena vacía. Un ejemplo
de cláusula es
\[
  A(XYZ, W) \to B(X)\, C(XY, Z)\, D(W).
\]

La regla anterior tiene el siguiente significado: el predicado $A$
recibe un vector de dimensión $2$; a partir de las cadenas del vector,
se asignan valores a las variables $X$, $Y$, $Z$ y $W$, y con esos
valores se construyen los vectores que reciben los predicados $B$, $C$
y $D$.

El proceso se ilustra con un ejemplo en el que el predicado $A$ recibe
el vector de cadenas $[abc, d]$. El primer paso es asignar las cadenas
del vector a los argumentos del predicado. El primer argumento de $A$
es $abc$ y el segundo es $d$. Como el primer argumento de $A$ está
definido como $XYZ$ y el segundo como $W$, debe cumplirse que $XYZ = abc$
y $W = d$. Esto significa que las variables $X$, $Y$ y $Z$ deben tomar
los valores posibles cuya concatenación sea $abc$, mientras que $W$ solo
puede tomar el valor $d$. En el cuadro~\ref{tab:valores-xyzw} se
muestran algunas asignaciones posibles para las variables $X$, $Y$, $Z$
y $W$ a partir del vector de entrada.

\begin{table}[ht]
  \centering
  \begin{tabular}{|c|c|c|c|}
    \hline
    $X$ & $Y$ & $Z$ & $W$ \\ \hline
    $\varepsilon$ & $\varepsilon$ & $abc$ & $d$ \\ \hline
    $\varepsilon$ & $a$ & $bc$ & $d$ \\ \hline
    $\varepsilon$ & $ab$ & $c$ & $d$ \\ \hline
    $\varepsilon$ & $abc$ & $\varepsilon$ & $d$ \\ \hline
    $a$ & $\varepsilon$ & $bc$ & $d$ \\ \hline
    $a$ & $b$ & $c$ & $d$ \\ \hline
    $a$ & $bc$ & $\varepsilon$ & $d$ \\ \hline
    $ab$ & $\varepsilon$ & $c$ & $d$ \\ \hline
    $ab$ & $c$ & $\varepsilon$ & $d$ \\ \hline
    $abc$ & $\varepsilon$ & $\varepsilon$ & $d$ \\ \hline
  \end{tabular}
  \caption{Asignaciones posibles para las variables $X$, $Y$, $Z$ y
  $W$ a partir del vector $[abc, d]$.}
  \label{tab:valores-xyzw}
\end{table}

Para continuar con el ejemplo, se supone que los valores de $X$, $Y$ y
$Z$ son $X = ab$, $Y = \varepsilon$ y $Z = c$. A partir de esta
asignación se construyen los vectores con los que se evalúan los
predicados del lado derecho $B(X)\, C(XY, Z)\, D(W)$. Como $X = ab$,
$Y = \varepsilon$, $Z = c$ y $W = d$, el lado derecho se evalúa con los
vectores $[ab]$, $[ab, c]$ y $[d]$, y el resultado es la derivación
$B(ab)\, C(ab, c)\, D(d)$.

El proceso de asignar los vectores a los argumentos, identificar los
valores de las variables e instanciar las partes derechas se repite en
cada uno de los predicados del lado derecho de la cláusula.

Existe otro tipo de producciones, denominadas cláusulas terminales, en
las que un predicado deriva en $\varepsilon$ para determinados valores
de sus argumentos. Por ejemplo,
\[
  B(ab) \to \varepsilon, \qquad
  C(ab, c) \to \varepsilon, \qquad
  D(d) \to \varepsilon.
\]
La forma general de estas producciones es
$A(X_1, \dots, X_n) \to \varepsilon$.

Un predicado reconoce un vector de cadenas si existe una asignación de
sus argumentos a valores de las variables y una secuencia de
derivaciones que termina en la cadena vacía. Por ejemplo, en el caso de
$B(ab) \to \varepsilon$, el predicado $B$ reconoce la cadena $ab$.

A modo de ilustración, se considera el siguiente conjunto de
producciones:
\begin{enumerate}
  \item $A(XYZ, W) \to B(X)\, C(XY, Z)\, D(W)$,
  \item $B(ab) \to \varepsilon$,
  \item $C(ab, c) \to \varepsilon$,
  \item $D(d) \to \varepsilon$.
\end{enumerate}
Con estas producciones, el predicado $A$ reconoce el vector $[abc, d]$:
con la asignación $X = ab$, $Y = \varepsilon$, $Z = c$ y $W = d$, el
predicado $A$ deriva en $B(ab)\, C(ab, c)\, D(d)$, y cada uno de los
predicados $B(ab)$, $C(ab, c)$ y $D(d)$ deriva en la cadena vacía.

Un predicado no reconoce un vector de cadenas cuando para ninguna
asignación de variables se deriva en la cadena vacía.

Presentadas estas nociones, a continuación se introducen las
definiciones formales de las gramáticas de concatenación de rango.

\subsection{Definiciones formales}

\begin{definicion}
Un \textbf{rango} es una tupla $(i, j)$ que representa un intervalo de
posiciones en una cadena, donde $i$ y $j$ son enteros no negativos tales
que $i \le j$.
\end{definicion}

Por ejemplo, con índices indexados en $0$, para la cadena $abcd$ el
rango $(1, 3)$ representa la subcadena $bc$.

\begin{definicion}
Una \textbf{gramática de concatenación de rango} se define como una
$5$-tupla
\[
  G = (N, T, V, P, S),
\]
donde:
\begin{itemize}
  \item $N$ es un conjunto finito de \textbf{predicados} o símbolos no
    terminales; cada predicado tiene una aridad, que indica la dimensión
    del vector de cadenas que reconoce;
  \item $T$ es un conjunto finito de \textbf{símbolos terminales};
  \item $V$ es un conjunto finito de \textbf{variables};
  \item $P$ es un conjunto finito de \textbf{cláusulas} de la forma
    \[
      A(x_1, x_2, \dots, x_k) \to B_1(y_{1,1}, y_{1,2}, \dots, y_{1,m_1})
      \dots B_n(y_{n,1}, y_{n,2}, \dots, y_{n,m_n}),
    \]
    donde $A, B_i \in N$, $x_i, y_{i,j} \in (V \cup T)^*$, y $k$ es la
    aridad de $A$;
  \item $S \in N$ es el \textbf{predicado inicial} de la gramática, que
    siempre tiene aridad $1$.
\end{itemize}
\end{definicion}

Por ejemplo, a continuación se muestra una gramática de concatenación
de rango que reconoce el lenguaje $L_{\mathrm{copy}}^3 = \{ w^3 \mid w
\in \{a, b, c\}^* \}$:
\[
  G_{\mathrm{copy}}^3 = (N, T, V, P, S),
\]
donde $N = \{A, S\}$, $T = \{a, b, c\}$, $V = \{X, Y, Z\}$, el símbolo
inicial es $S$ y el conjunto de cláusulas $P$ es el siguiente:
\begin{enumerate}
  \item $S(XYZ) \to A(X, Y, Z)$,
  \item $A(aX, aY, aZ) \to A(X, Y, Z)$,
  \item $A(bX, bY, bZ) \to A(X, Y, Z)$,
  \item $A(cX, cY, cZ) \to A(X, Y, Z)$,
  \item $A(\varepsilon, \varepsilon, \varepsilon) \to \varepsilon$.
\end{enumerate}

A diferencia de las gramáticas presentadas en el capítulo anterior, que
se definen a partir de un mecanismo de generación de cadenas, las RCG
se formulan directamente como un mecanismo de reconocimiento. Su
funcionamiento consiste en determinar si una cadena pertenece o no al
lenguaje.

\begin{definicion}
Una \textbf{sustitución de rango} es un mecanismo que reemplaza una
variable por un rango de la cadena, respetando la estructura del
argumento que se asocia a la cadena que se reconoce.
\end{definicion}

Por ejemplo, dado el predicado $A(Xa)$ con $X \in V$ y $a \in T$, la
estructura del argumento de $A$ es una variable $X$ seguida del
terminal $a$. Si el predicado $A$ recibe la cadena $baa$, entonces $X$
se puede asociar con la subcadena $ba$ de la cadena original: si
$X = ba$, entonces $Xa = baa$. La variable $X$ no puede tomar el valor
$baa$, porque ningún carácter de la cadena de entrada coincide con el
terminal $a$ al final del argumento. La variable $X$ tampoco puede
tomar el valor $b$, porque con ese valor el argumento $Xa$ no cubre la
cadena completa.

\begin{definicion}
Las \textbf{gramáticas de concatenación de rango simple} (SRCG) son un
subconjunto de las RCG que restringen la forma de las cláusulas. Una
RCG $G$ es \textbf{simple} si los argumentos del lado derecho de cada
cláusula son variables distintas, y todas estas variables aparecen una
sola vez en los argumentos del lado izquierdo.
\end{definicion}

Las SRCG son equivalentes a los sistemas lineales de reescritura libres
del contexto (LCFRS) introducidos en~\cite{vijayshanker1987lcfrs}. Este
subconjunto se utiliza en~\cite{aguilera2016srcg} para describir el
orden de las variables de una fórmula booleana en una variante del SAT.

Una vez introducidas las definiciones formales, a continuación se
presenta un ejemplo completo de reconocimiento de una cadena por una
RCG, aplicando los mecanismos descritos en esta sección.

\subsection{Ejemplo de reconocimiento}

La presentación intuitiva describió cómo se asocian los elementos de
un vector a los argumentos de un predicado y cómo se construyen los
vectores que reciben los predicados del lado derecho de una cláusula.
En esta subsección se aplica ese proceso a un caso completo: el
reconocimiento de la cadena $abcabcabc$ por la gramática
$G_{\mathrm{copy}}^3$ presentada anteriormente.

La cadena $abcabcabc$ se reconoce por $G_{\mathrm{copy}}^3$ mediante
la siguiente secuencia de derivaciones:
\[
  S(abcabcabc) \to A(abc, abc, abc) \to A(bc, bc, bc) \to A(c, c, c)
  \to A(\varepsilon, \varepsilon, \varepsilon) \to \varepsilon.
\]
A continuación se describen los pasos.
\begin{itemize}
  \item \textbf{Primer paso:} en la primera cláusula
    $S(XYZ) \to A(X, Y, Z)$, la sustitución de rango asocia las
    variables $X$, $Y$, $Z$ a los valores $X = abc$, $Y = abc$ y
    $Z = abc$. De esta forma se deriva en el predicado $A(abc, abc,
    abc)$.
  \item \textbf{Segundo paso:} en la segunda cláusula
    $A(aX, aY, aZ) \to A(X, Y, Z)$, la sustitución de rango asocia las
    variables $X$, $Y$, $Z$ a los valores $X = bc$, $Y = bc$ y
    $Z = bc$. Con estos valores se deriva en el predicado
    $A(bc, bc, bc)$.
  \item \textbf{Tercer paso:} en la tercera cláusula
    $A(bX, bY, bZ) \to A(X, Y, Z)$, la sustitución de rango asocia las
    variables $X$, $Y$, $Z$ a los valores $X = c$, $Y = c$ y $Z = c$.
    Con estos valores se deriva en el predicado $A(c, c, c)$.
  \item \textbf{Cuarto paso:} en la cuarta cláusula
    $A(cX, cY, cZ) \to A(X, Y, Z)$, la sustitución de rango asocia las
    variables $X$, $Y$, $Z$ a los valores $X = \varepsilon$,
    $Y = \varepsilon$ y $Z = \varepsilon$. Con estos valores se deriva
    en el predicado $A(\varepsilon, \varepsilon, \varepsilon)$.
  \item \textbf{Quinto paso:} en la última cláusula
    $A(\varepsilon, \varepsilon, \varepsilon) \to \varepsilon$ se deriva
    en la cadena vacía, con lo que la cadena $abcabcabc$ queda
    reconocida.
\end{itemize}

Una vez presentado el proceso de derivación, en la siguiente sección se
analizan las propiedades del formalismo: clausura bajo operaciones sobre
conjuntos, poder expresivo y complejidad del problema de la palabra.

\section{Propiedades de las RCG}
\label{sec:propiedades-rcg}

Las RCG fueron diseñadas como un formalismo modular, en el sentido de
que las operaciones de unión, intersección y complemento son cerradas
para este formalismo~\cite{boullier1998proposal}. La ausencia de esta
modularidad en las gramáticas libres del contexto fue una de las
limitaciones que motivó la propuesta original de las
RCG~\cite{boullier1998proposal}. En esta sección se describen las
propiedades de clausura que demuestran la modularidad del formalismo y
la caracterización en términos de complejidad.

\begin{teorema}[Clausura bajo operaciones de conjuntos]
Las RCG son cerradas bajo unión, intersección y complemento. La unión
y la intersección de dos RCG da como resultado un formalismo que
pertenece a las RCG, mientras que el formalismo resultante del
complemento de una RCG es también una RCG~\cite{boullier1998proposal}.
\end{teorema}

En cuanto al poder expresivo, las RCG son estrictamente más expresivas
que las gramáticas libres del contexto. Por ejemplo, para todo
$k \ge 2$, el lenguaje $L_{\mathrm{copy}}^k = \{ w^k \mid w \in \Sigma^*\}$
es dependiente del contexto y no libre del
contexto~\cite{hopcroft2006introduction}, y se reconoce por una RCG. En
particular, el lenguaje $L_{\mathrm{copy}}^3$ se reconoce por la
gramática $G_{\mathrm{copy}}^3$ del ejemplo anterior.

Una vez presentadas las propiedades estructurales, a continuación se
analiza la caracterización del formalismo en términos de complejidad
computacional, que constituye el resultado que motiva el uso de las RCG
en este trabajo.

\subsection{Complejidad y caracterización en términos de $PTIME$}

El siguiente resultado es el pilar que conecta a las RCG con la teoría
de la complejidad y sobre el que se apoyan los resultados de los
capítulos posteriores.

\begin{teorema}[Caracterización en términos de $PTIME$]
\label{teo:prcg-ptime}
Un lenguaje $L$ se reconoce por una RCG en su variante positiva si y
solo si $L$ admite un algoritmo de reconocimiento que se ejecuta en
tiempo polinomial~\cite{bertsch2001rcg, boullier2000range}.
\end{teorema}

La demostración del teorema~\ref{teo:prcg-ptime} se realiza
en~\cite{bertsch2001rcg, boullier2000range}. La dirección que afirma
que todo lenguaje RCG está en $PTIME$ se sigue de la existencia de un
algoritmo de reconocimiento polinomial para el formalismo, que se
describe a continuación.

El algoritmo de reconocimiento estándar para las RCG utiliza un proceso
de memorización sobre las cadenas asignadas a los argumentos de los
predicados~\cite{boullier1998proposal}. La cantidad máxima de rangos de
una cadena de longitud $n$ es $n^2$, y la máxima aridad de un predicado
es constante, por lo que el proceso de memorización emplea una cantidad
polinomial de estados. La complejidad del algoritmo es $O(|P|
n^{2h(l+1)})$, donde $h$ es la máxima aridad de un predicado, $l$ es
la máxima cantidad de predicados en el lado derecho de una cláusula y
$n$ es la longitud de la cadena que se reconoce. El teorema establece,
además, la dirección inversa: todo lenguaje en $PTIME$ admite una RCG
que lo reconoce.

El teorema~\ref{teo:prcg-ptime} es el resultado que motiva la elección
de las RCG como formalismo de base para este trabajo. El alcance exacto
de la caracterización y sus consecuencias para la relación entre
formalismos gramaticales y problemas de la clase $NP$ son el objeto
central de los capítulos posteriores.

Una vez establecida la caracterización de las RCG como formalismo que
reconoce exactamente la clase $PTIME$, en la siguiente sección se
revisan los antecedentes directos del trabajo.

\section{Antecedentes en el estudio del SAT mediante formalismos de la teoría de lenguajes}
\label{sec:antecedentes}

El estudio del problema de la satisfacibilidad booleana mediante
formalismos de la teoría de lenguajes cuenta con antecedentes directos
en la Facultad de Matemática y Computación de la Universidad de La
Habana. En esta sección se revisan los tres trabajos previos que
constituyen el punto de partida del presente trabajo.

En~\cite{fernandez2007sat} se propone una estrategia para resolver
instancias del SAT mediante gramáticas libres del contexto. La
estrategia consiste en tres pasos: asumir que todas las variables de la
fórmula son distintas; construir un autómata finito que reconozca las
cadenas de $0$ y $1$ que hacen verdadera la fórmula bajo ese supuesto;
e intersectar el lenguaje regular obtenido con un lenguaje libre del
contexto que obligue a que todas las instancias de una misma variable
tomen el mismo valor. La intersección resultante es un lenguaje libre
del contexto sobre el que el problema del vacío se decide en tiempo
polinomial, con lo que el procedimiento completo opera en complejidad
$O(n^3)$ para ciertas subclases del SAT.

En~\cite{aguilera2016srcg} se extiende la idea anterior al emplear
gramáticas de concatenación de rango simple en lugar de gramáticas
libres del contexto. El autómata booleano asociado a la fórmula se
intersecta con una SRCG que modela el orden de las variables, lo que
amplía el conjunto de fórmulas booleanas que admiten solución mediante
esta estrategia.

En~\cite{gomez2025sat} se propone una estrategia diferente: en lugar de
resolver el problema del vacío para el lenguaje de las interpretaciones
que hacen verdadera a una fórmula dada, se construye el lenguaje de
todas las fórmulas booleanas satisfacibles y se decide la
satisfacibilidad mediante el problema de la palabra. Como parte de
dicho trabajo se propone una gramática que reconocería el lenguaje de
las fórmulas satisfacibles, y a partir de ella se derivan afirmaciones
sobre la expresividad del formalismo con respecto a la clase $NP$.

Los tres trabajos anteriores conforman un programa de investigación
acumulativo en la Facultad sobre la relación entre formalismos
gramaticales y el SAT. El presente trabajo se inscribe en este programa
y se propone examinar con precisión formal la clasificación de la
gramática propuesta en~\cite{gomez2025sat} dentro de la jerarquía de
formalismos considerada en este capítulo, analizar las consecuencias de
dicha clasificación para los resultados derivados, y explorar
construcciones de RCG para variantes polinomiales del SAT. Estos
objetivos se desarrollan en los capítulos siguientes.

\bibliographystyle{plain}
\bibliography{referencias}

\end{document}
